\chapter{Market Microstructure and Stochastic Foundations of Liquidity Provision}
\label{ch:microstructure}

\section{Introduction}
High-frequency market making is fundamentally an exercise in stochastic optimal control under strict inventory constraints. This chapter establishes the quantitative framework required to model the dynamics of the Limit Order Book (LOB) and derives the theoretical foundations for optimal liquidity provision. We formulate the market-making problem as a utility maximization task, focusing on the derivation of optimal bid and ask offsets from the fair value of an asset.

\section{Dynamics of the Limit Order Book}
The LOB is modeled as a continuous-time double auction where liquidity is provided via limit orders and consumed by market orders. Let $L_t$ denote the state of the LOB at time $t$. The evolution of the LOB is governed by the arrival of orders following a point process.

\subsection{Price-Time Priority and Execution Mechanics}
Execution at major electronic venues is governed by Price-Time Priority. An order $O(p, q, \tau)$ with price $p$, quantity $q$, and arrival time $\tau$ is executed only if:
\begin{enumerate}
    \item There exists an aggressive market order at price $p'$ such that $p' \geq p$ (for sell orders) or $p' \leq p$ (for buy orders).
    \item No other order $O'$ exists with price $p$ and arrival time $\tau' < \tau$.
\end{enumerate}
This deterministic rule underscores the necessity of the ultra-low-latency execution implemented in this research. Queue position is a primary determinant of fill probability; therefore, any delay in quoting response significantly increases the risk of adverse selection and missed execution alpha.

\section{The Avellaneda--Stoikov Optimal Control Framework}
The core financial engineering problem is to determine the optimal bid-ask spread that maximizes the market maker's terminal utility. We build upon the seminal work of \citet{Avellaneda2008}.

\subsection{Asset Price and Wealth Dynamics}
Let $S_t$ be the mid-price of an asset, following a Wiener Process:
\begin{equation}
    dS_t = \sigma dW_t
\end{equation}
where $\sigma$ denotes the instantaneous volatility and $W_t$ is a standard Brownian motion. The market maker's wealth $X_t$ and inventory $q_t$ evolve according to the following stochastic differential equations (SDEs):
\begin{equation}
    dX_t = (S_t + \delta^a_t) dN^a_t - (S_t - \delta^b_t) dN^b_t
\end{equation}
\begin{equation}
    dq_t = dN^b_t - dN^a_t
\end{equation}
where $\delta^a_t, \delta^b_t$ are the offsets from the mid-price, and $N^a_t, N^b_t$ are Poisson processes representing the arrival of sell and buy fill events, respectively.

\subsection{Utility Maximization and the HJB Equation}
The market maker seeks to maximize the expected exponential utility of terminal wealth:
\begin{equation}
    \max_{\delta^a, \delta^b} \mathbf{E} \left[ -\exp(-\gamma (X_T + q_T S_T)) \right]
\end{equation}
where $\gamma$ is the risk-aversion parameter. The value function $u(s, x, q, t)$ satisfies the Hamilton-Jacobi-Bellman (HJB) equation:
\begin{equation}
    \frac{\partial u}{\partial t} + \frac{1}{2} \sigma^2 \frac{\partial^2 u}{\partial s^2} + \max_{\delta^b} \lambda^b(\delta^b) [u(s, x-s+\delta^b, q+1, t) - u] + \max_{\delta^a} \lambda^a(\delta^a) [u(s, x+s+\delta^a, q-1, t) - u] = 0
\end{equation}
where $\lambda(\delta)$ is the fill intensity function, modeled as a decreasing function of the offset:
\begin{equation}
    \lambda(\delta) = A e^{-k\delta}
\end{equation}
with $A$ and $k$ representing market liquidity parameters.

\subsection{The Reservation Price and Optimal Offsets}
The \textbf{Reservation Price} $r(s, q, t)$ is the price at which the market maker is indifferent between their current portfolio and a portfolio with one fewer unit of inventory. Under the assumption of constant volatility and small $\gamma$, the reservation price is approximated as:
\begin{equation}
    r(s, q, t) = S_t - q \gamma \sigma^2 (T - t)
\end{equation}
This derivation reveals the fundamental mechanics of inventory management: as $q_t$ increases, the reservation price shifts downward, inducing the market maker to lower their bid and ask quotes to encourage sells and discourage further buys. The optimal symmetric spread $\psi^*$ around the reservation price is given by:
\begin{equation}
    \psi^* = \frac{2}{\gamma} \ln \left( 1 + \frac{\gamma}{k} \right) + \gamma \sigma^2 (T - t)
\end{equation}

\section{Microstructural Signals and Adverse Selection}
Classical models assume that order arrivals are independent of the LOB state. However, in modern markets, the probability of a fill is highly contingent on microstructural signals.

\subsection{Order Book Imbalance (OBI)}
The Order Book Imbalance $\rho_t$ is a primary signal used in this research to proxy information asymmetry:
\begin{equation}
    \rho_t = \frac{V_t^b - V_t^a}{V_t^b + V_t^a}
\end{equation}
where $V^b_t, V^a_t$ are the volumes at the best bid and ask. A significant imbalance ($\rho_t \rightarrow 1$ or $\rho_t \rightarrow -1$) indicates a directional bias in the aggressive order flow, suggesting an imminent price innovation.

\subsection{Adverse Selection and Toxic Flow}
Adverse selection occurs when a market maker provides liquidity to an informed trader who possesses a short-term predictive edge. This "toxic flow" results in the market maker being filled just before a large price move against their position. To mitigate this risk, the adaptive reinforcement learning model (detailed in Chapter \ref{ch:rl_strategy}) utilizes $\rho_t$ and $\sigma_t$ as non-linear features to adjust quoting spreads dynamically, surpassing the performance of static AS-heuristics.

\section{Conclusion}
This chapter has established the stochastic foundations of optimal liquidity provision. The derivation of the reservation price and optimal spread provides the theoretical benchmark for our system. In the following chapters, we detail how these quantitative models are operationalized using high-fidelity data generation and hardware-accelerated reinforcement learning.
